% -----------------------------------------------------------
\section{Tempt, Temptation}
% -----------------------------------------------------------
	\label{TEMPTATION}
	
% ----------------------
\subsection{See also}
% ----------------------

	\hyperref[ACCOUNTABILITY]{Accountability}, \hyperref[AGENCY]{Agency},

% ----------------------
\subsection{1828 American Dictionary of the English Language} % *1828
% ----------------------
	\label{TEMPTATION-1828}
	
	% -----
	\subsubsection{Tempt}
	% -----

		\begin{compactitem}
			\item To incite or solicit to an evil act; to entice to something wrong by presenting arguments that are plausible or convincing, or by the offer of some pleasure or apparent advantage as the inducement.
			\item To provoke; to incite.
			\item To solicit; to draw; without the notion of evil.
			\item To try; to venture on; to attempt.
			\item In Scripture, to try; to prove; to put to trial for proof.
		\end{compactitem}

% % ----------------------
% \subsection{Oxford English Dictionary} % *OED
% % ----------------------
	% \label{TEMPTATION-OED}

	% \begin{compactitem}
		% \item[]
	% \end{compactitem}
	
% % ----------------------
% \subsection{Restoration Lexicon} % *RL *RESTORATION LEXICON
% % ----------------------
	% \label{TEMPTATION-OED}

	% \begin{compactitem}
		% \item[]
	% \end{compactitem}

% ----------------------
\subsection{Scriptural Usage} % *SCRIPTURES
% ----------------------
	\label{TEMPTATION-SCRIPTURES}

	% % -----
	% \subsubsection{Old Testament} % *OT
	% % -----
		% \label{TEMPTATION-OT}
	
		% % --
		% \paragraph{KJV}
		% % --
			% \label{TEMPTATION-OT-KJV}
	
			% \begin{tabular}{ll}
				% Total occurrences in the OT & 
			% \end{tabular}
			
			% \begin{compactdesc}
				% \item[]
			% \end{compactdesc}
			
		% --
		\paragraph{Hebrew Bible}
		% --
			\label{TEMPTATION-HEBREW}
		
			%
			\subparagraph{\texorpdfstring{\he{nAsAh}}{}} % paste actual hebrew text here
			%
				\label{NASAH} % whatever is in step bible without periods
			
				\begin{compactdesc}
					\item[HALOT 702a, PDF 758] \textbf{} % Use the 1994 PDF
						\begin{compactitem}
							\item[1] to \textbf{put someone to the test}
							\item[2] to give experience, \textbf{train}
							\item[3] to \textbf{conduct a test, make an attempt}
						\end{compactitem}
					% \item[BDB , PDF ] \textbf{}
						% \begin{compactitem}
							% \item[1]
						% \end{compactitem}
					% \item[DCH PDF ] \textbf{} % don't include the actual page number, they repeat from volume to volume
						% \begin{compactitem}
							% \item[1]
						% \end{compactitem}
					\item[TDOT 9:443, PDF 465]
				\end{compactdesc}
				
				\begin{compactdesc}
					\item[Some OT uses] \textbf{}
						\begin{compactdesc}
							\item[{\hyperref[GENESIS-22-1]{Genesis 22:1}}]
						\end{compactdesc}
				\end{compactdesc}
			
		% % --
		% \paragraph{Septuagint}
		% % --
			% \label{TEMPTATION-LXX}
		
			% %
			% \subparagraph{\texorpdfstring{\gr{sample word}}{}}
			% %
		
	% -----
	\subsubsection{New Testament} % *NT
	% -----
		\label{TEMPTATION-NT}
	
		% % --
		% \paragraph{KJV}
		% % --
			% \label{TEMPTATION-NT-KJV}		
	
			% \begin{tabular}{ll}
				% Total occurrences in the NT & 
			% \end{tabular}
			
			% \begin{compactdesc}
				% \item[]	
			% \end{compactdesc}
			
		% --
		\paragraph{Greek New Testament} % *GREEK
		% --
			\label{TEMPTATION-GREEK}
			
			%
			\subparagraph{\texorpdfstring{\gr{ἐκπειράζω}}{}}
			%
				\label{EKPEIRAZO}
			
				\begin{compactdesc}
					\item[BDAG 704a] \textbf{}
						\begin{compactitem}
							\item[1] \textbf{to subject to test or proof, \textit{tempt}}
							\item[2] \textbf{to entrap someone into giving information that will jeopardize the person, \textit{entrap}}
							\item[3] \textbf{to entice to do wrong by offering attractive benefits, \textit{tempt}}
						\end{compactitem}
					% \item[CGL , PDF ] \textbf{}
						% \begin{compactitem}
							% \item 
						% \end{compactitem}
					% \item[LSJ , PDF ] \textbf{}
						% \begin{compactitem}
							% \item 
						% \end{compactitem}
				\end{compactdesc}
				
				% \begin{tabular}{ll}
					% Total occurrences in the NT & 
				% \end{tabular}
				
				\begin{compactdesc}
					\item[Some NT uses] \textbf{}
						\begin{compactdesc}
							\item[{\hyperref[MATTHEW-4-1]{Matthew 4:1}}]
							\item[{\hyperref[LUKE-4-2]{Luke 4:2}}]
						\end{compactdesc}
				\end{compactdesc}
		
			%
			\subparagraph{\texorpdfstring{\gr{πειράζω}}{}}
			%
				\label{PEIRAZO}
			
				\begin{compactdesc}
					\item[BDAG 704a] \textbf{}
						\begin{compactitem}
							\item[1] \textbf{to make an effort to do something, \textit{try, attempt}} at times in a context indicating futility
							\item[2] \textbf{to endeavor to discover the nature or character of something by testing, \textit{try, make trial of, put to the test}}
							\item[2.a] general \gr{τινά} \textit{someone}
							\item[2.b] of God or Christ, who put people to the test, in a favorable sense...so that they may prove themselves true...also of painful trials sent by God
							\item[2.c] The Bible...also speaks of a trial of God by humans. Their intent is to put God to the test, to discover whether God really can do a certain thing, especially whether God notices sin and is able to punish it
								\begin{compactitem}
									\item I thought that the LXX of \hyperref[MALACHI-3-10]{Malachi 3:10} might use this word, but there it's a form of \gr{ἐπισκέπτομαι}. BSV has \textit{probo} though.
								\end{compactitem}
							\item[3] \textbf{to attempt to entrap through a process of inquiry, \textit{test}}. Jesus was so treated by his opponents, who planned to use their findings against him
							\item[4] \textbf{to entice to improper behavior, \textit{tempt}}...above all the devil works in this way; hence he is directly called \gr{ὁ πειράζων} \textit{the tempter}...he tempts humans...but he also makes bold to tempt Jesus
								\begin{itemize}
									\item \hyperref[MATTHEW-4-1]{Matthew 4:1} is listed here.
								\end{itemize}
						\end{compactitem}
					% \item[CGL , PDF ] \textbf{}
						% \begin{compactitem}
							% \item 
						% \end{compactitem}
					\item[LSJ 1354b, PDF 1425] \textbf{}
						\begin{compactitem}
							\item[I.1] \textit{make proof} or \textit{trial of}
							\item[I.2] with infinitive, \textit{attempt} to do
							\item[I.3] passive...\textit{are tried, proved}...\textit{to be experienced}
							\item[II.1] with accusative of person, \textit{try, tempt} a person, \textit{put} him \textit{to the test}
							\item[II.2] in bad sense, \textit{seek to seduce, tempt}...passive, \textit{to be sorely tried}...\textit{to be tempted to sin}
						\end{compactitem}
				\end{compactdesc}
				
				% \begin{tabular}{ll}
					% Total occurrences in the NT & 
				% \end{tabular}
				
				\begin{compactdesc}
					\item[Some NT uses] \textbf{}
						\begin{compactdesc}
							\item[{\hyperref[MATTHEW-4-1]{Matthew 4:1}}]
						\end{compactdesc}
				\end{compactdesc}
		
	% -----
	\subsubsection{Book of Mormon} % *BOM
	% -----
		\label{TEMPTATION-BOM}
	
		% \begin{tabular}{ll}
			% Total occurrences in the BOM & 
		% \end{tabular}
		
		\begin{compactdesc}
			\item[{\hyperref[1NEPHI-12-16]{1 Nephi 12:16}}]
			\item[{\hyperref[3NEPHI-6-17]{3 Nephi 6:17}}]
			\item[{\hyperref[3NEPHI-28-39]{3 Nephi 28:39}}]
		\end{compactdesc}
		
	% -----
	\subsubsection{Doctrine and Covenants} % *DC
	% -----
		\label{TEMPTATION-DC}
	
		% \begin{tabular}{ll}
			% Total occurrences in the DC & 
		% \end{tabular}
		
		\begin{compactdesc}
			\item[{\hyperref[DC101-28]{DC 101:28}}]
		\end{compactdesc}
		
	% % -----
	% \subsubsection{Pearl of Great Price} % *PGP
	% % -----
		% \label{TEMPTATION-PGP}
	
		% \begin{tabular}{ll}
			% Total occurrences in the PGP & 
		% \end{tabular}
		
		% \begin{compactdesc}
			% \item[]
		% \end{compactdesc}
		
% % ----------------------
% \subsection{Key Phrases} % *KEYPHRASES *PHRASES
% % ----------------------
	% \label{TEMPTATION-PHRASES}

	% \begin{compactdesc}
		% \item[] \textbf{\#times} | list
			% % \begin{compactdesc}
				% % \item[Bible anchor verses] \phantom{.}
					% % \begin{compactdesc}
						% % \item[Romans 5:5] \gr{}
					% % \end{compactdesc}
				% % \item[Commentary] ---
			% % \end{compactdesc}
	% \end{compactdesc}
	
% % ----------------------
% \subsection{Bible Dictionaries} % *DICTIONARIES
% % ----------------------
	% \label{TEMPTATION-BD}

% % ----------------------
% \subsection{Restoration Usage} % *RESTORATION
% % ----------------------
	% \label{TEMPTATION-RESTORATION}

	% % -----
	% \subsubsection{Joseph Smith} % *JS
	% % -----
		% \label{TEMPTATION-JS}
	
		% \begin{compactdesc}
			% \item[{\hyperref[KEY]{Date/Source}}]
		% \end{compactdesc}
		
	% % -----
	% \subsubsection{Temple Ordinances}
	% % -----
		% \label{TEMPTATION-TEMPLE}
	
		% \begin{compactdesc}
			% \item[{\hyperref[KEY]{Date/Source}}]
		% \end{compactdesc}
	
	% % -----
	% \subsubsection{Brigham Young} % *BY
	% % -----
		% \label{TEMPTATION-BY}
	
		% \begin{compactdesc}
			% \item[{\hyperref[KEY]{Date/Source}}]
		% \end{compactdesc}
	
% ----------------------
\subsection{Commentary and Studies} % *COMMENTARY *STUDIES
% ----------------------
	\label{TEMPTATION-COMMENTARY}

	% -----
	\subsubsection{Key Points}
	% -----
		\label{TEMPTATION-KEYS}
	
		\begin{compactitem}
			\item Satan does not have ``power'' to tempt little children; see \hyperref[DC29-46]{DC 29:46--47} and \hyperref[MOSES-6-55]{Moses 6:55}.
		\end{compactitem}
		
	% -----
	\subsubsection{Christ's experience with temptation}
	% -----
		\label{TEMPTATION-christ}
		
		\begin{compactitem}
			\item \hyperref[HEBREWS-4-14]{Hebrews 4:14--15}
			\item \hyperref[MOSIAH-3-7]{Mosiah 3:7}
			\item \hyperref[MOSIAH-15-5]{Mosiah 15:5}
			\item \hyperref[ALMA-7-11]{Alma 7:11}
			\item \hyperref[DC20-22]{DC 20:22}
		\end{compactitem}